\section{Structure-Preserving Semismooth Newton Method}\label{sec:ssn}

\subsection{Nonsmooth residual and generalized Jacobian}

Let $H:\RR^n\rightarrow\RR^n$ denote the complete residual in
\eqref{eq:complete-system}. Smooth ac/dc network equations and converter loss
functions are combined with the locally Lipschitz Fischer--Burmeister and
median maps. The resulting $H$ is differentiable away from switching ties and
semismooth at the ties under the standard NCP construction
\cite{facchinei2003complementarity,qi1993nonsmooth}.

At iteration $k$, one element
\begin{equation}
 J_k\in\Bsub H(w_k)
 \label{eq:selected-jacobian}
\end{equation}
is assembled analytically. The network part follows the production convention
$J_{\mathrm{net}}=\partial(\mathrm{calculated}-\mathrm{specified})/\partial x$
with right-hand side equal to the network mismatch. Local converter rows store
the conventional derivative of their zero residual and use right-hand side
$-H_c$. This sign split is resolved during assembly into one linear system
\begin{equation}
 J_k d_k=-H(w_k).
 \label{eq:newton-step}
\end{equation}

The exact Fischer--Burmeister map is not regularized by a hidden epsilon. At
the biactive origin, the implementation selects the symmetric Clarke element
\begin{equation}
 \left(\frac{1}{\sqrt 2}-1,
       \frac{1}{\sqrt 2}-1\right)\in
 \Clarke\phi_{\mathrm{FB}}(0,0).
 \label{eq:fb-origin}
\end{equation}
For the median map, the active branch determines the derivative; equal
branches use an equal-weight finite selection. This policy makes the selected
matrix deterministic while retaining the fixed equation layout.

\subsection{Optional smooth continuation}

Degenerate intersections may cause an abrupt change in the selected
generalized Jacobian. The optional continuation replaces
\eqref{eq:fb} with the Kanzow smoothing
\begin{equation}
 \phi_{\mu}(a,b)
 =\sqrt{a^2+b^2+2\mu^2}-a-b,
 \qquad \mu>0,
 \label{eq:smooth-fb}
\end{equation}
and replaces nested min/max operations in \eqref{eq:median} with the
Chen--Harker--Kanzow--Smale smooth composition
\cite{kanzow1996continuation}. The parameter is reduced from $\mu_0$ to
$\mu_{\min}$ using a coarse/fine multiplicative schedule. After each reduction,
the residual and Jacobian are reassembled at the unchanged state. Only matrix
values change; variable order and sparse coordinates do not.

The reported operating point is always re-evaluated with $\mu=0$. Thus a small
smooth surrogate residual is not reported as an exact complementarity
certificate. Smooth continuation is a globalization aid, not a change to the
target physical model.

\subsection{Globalized iteration}

The production method combines the local semismooth step with residual and
variable scaling, a nonmonotone line search, and fail-closed fallbacks. Let
$D_F$ and $D_w$ denote residual and variable scaling matrices. The scaled
linear system is
\begin{equation}
 (D_FJ_kD_w^{-1})\widehat d_k=-D_FH(w_k),\qquad
 d_k=D_w^{-1}\widehat d_k.
 \label{eq:scaled-newton}
\end{equation}
The line search accepts $w_{k+1}=w_k+\alpha_k d_k$ when the merit function
$\Psi(w)=\tfrac12\|D_FH(w)\|_2^2$ satisfies a nonmonotone sufficient-decrease
test. Trust-region, pseudo-transient, active-set-seeded, and homotopy paths
remain external global safeguards; none changes the fixed local VSC state.

\begin{algorithm}[t]
\caption{Structure-preserving semismooth Newton solve}
\label{alg:ssn}
\KwIn{Hybrid system data, initial state $w_0$, tolerances, $\mu_0$}
Build the network and six-state-per-VSC sparse pattern once\;
Analyze the full and admissible reduced patterns once\;
\For{$k=0,1,\ldots,k_{\max}-1$}{
  Evaluate exact or $\mu_k$-smoothed $H(w_k)$\;
  \If{the exact target residual satisfies the stopping test}{return certified root\;}
  Assemble one selected analytic generalized Jacobian $J_k$\;
  Attempt the certified local-block Schur step\;
  \If{a Schur certificate fails}{solve the complete sparse system\;}
  Apply scaling and the globalized step acceptance rule\;
  \If{the smooth residual has decreased sufficiently and $\mu_k>\mu_{\min}$}{
    reduce $\mu_k$ and reassemble at the accepted state\;
  }
}
Return an explicit nonconvergence or infeasibility diagnosis\;
\end{algorithm}

\subsection{Structure preservation}

Each admitted converter always contributes six rows and six columns. A limit
activation, priority tie, or droop saturation modifies residual and Jacobian
values but not coordinates. Therefore the matrix memory layout, device-to-row
mapping, symbolic analysis, and result index remain stable. This property
removes repeated structural allocation and is also the prerequisite for the
local Schur method in Section~\ref{sec:schur}.

The fixed pattern does not imply that every numerical factorization can reuse
the same pivots. Numeric-only KLU refactorization is attempted only after a
previous factorization and is guarded by a normwise backward-error test. A
failed refactorization or an inaccurate solution triggers fresh pivoting.
